\usepackage[margin=0.85in,headheight=14pt]{geometry}
\usepackage[T1]{fontenc}
\usepackage[utf8]{inputenc}
\usepackage{lmodern}
\usepackage{amsmath}
\usepackage{amssymb}
\usepackage{mathtools}
\usepackage{siunitx}
\usepackage{booktabs}
\usepackage{tabularx}
\usepackage{longtable}
\usepackage{array}
\usepackage{enumitem}
\usepackage{microtype}
\usepackage{graphicx}
\usepackage[dvipsnames]{xcolor}
\usepackage{tikz}
\usepackage[american currents,american voltages]{circuitikz}
\usepackage{pgfplots}
\usepackage[most]{tcolorbox}
\usepackage{multicol}
\usepackage{fancyhdr}
\usepackage{caption}
\usepackage{subcaption}
\usepackage{placeins}
\usepackage[hidelinks]{hyperref}
\usepackage[capitalise,noabbrev]{cleveref}

\pgfplotsset{compat=1.18}
\usetikzlibrary{arrows.meta,calc,positioning}
\graphicspath{{figures/}}
\sisetup{per-mode=symbol,detect-all,retain-explicit-plus}

\definecolor{GuideBlue}{HTML}{244B5A}
\definecolor{GuideGray}{HTML}{F1F3F4}
\definecolor{GuideLine}{HTML}{7B858B}
\definecolor{GuideGreen}{HTML}{2F5D50}
\definecolor{GuideAmber}{HTML}{7A5C1E}
\definecolor{GuideRed}{HTML}{8C2D2D}
\definecolor{GuideTeal}{HTML}{1F6B6B}
\definecolor{GuideIndigo}{HTML}{433D78}
\definecolor{GuideSlate}{HTML}{51606B}
\definecolor{GuidePlum}{HTML}{6B3A5B}

\setlength{\parindent}{0pt}
\setlength{\parskip}{0.58em}
\setlist{nosep,leftmargin=*}
\setlist[description]{style=nextline,leftmargin=1.25in,labelwidth=1.15in}
\setcounter{secnumdepth}{2}
\setcounter{tocdepth}{2}
\raggedbottom
\emergencystretch=3em

\pagestyle{fancy}
\fancyhf{}
\lhead{\small Circuit Theory for the Amateur Extra Exam}
\rhead{\small \thepage}
\renewcommand{\headrulewidth}{0.35pt}

\captionsetup{font=small,labelfont=bf}

\tcbset{
  enhanced,
  sharp corners,
  boxrule=0.45pt,
  colback=GuideGray,
  colframe=GuideLine,
  left=7pt,
  right=7pt,
  top=6pt,
  bottom=6pt,
  before skip=0.75em,
  after skip=0.75em,
  breakable
}

% Callout boxes are colour-coded by type: each has its own accent frame and a
% faint matching tint, with a coloured title bar for quick visual scanning.
\newtcolorbox{controlsbox}{title=\textbf{Controls Connection},
  colframe=GuideBlue,   colback=GuideBlue!5,   coltitle=white, colbacktitle=GuideBlue}
\newtcolorbox{physicalbox}{title=\textbf{Physical Insight},
  colframe=GuideTeal,   colback=GuideTeal!5,   coltitle=white, colbacktitle=GuideTeal}
\newtcolorbox{energybox}{title=\textbf{Energy Insight},
  colframe=GuideGreen,  colback=GuideGreen!5,  coltitle=white, colbacktitle=GuideGreen}
\newtcolorbox{unitsbox}{title=\textbf{Units Check},
  colframe=GuideAmber,  colback=GuideAmber!6,  coltitle=white, colbacktitle=GuideAmber}
\newtcolorbox{exambox}{title=\textbf{Exam Relevance},
  colframe=GuideIndigo, colback=GuideIndigo!5, coltitle=white, colbacktitle=GuideIndigo}
\newtcolorbox{mistakebox}{title=\textbf{Common Mistake},
  colframe=GuideRed,    colback=GuideRed!5,    coltitle=white, colbacktitle=GuideRed}
\newtcolorbox{advancedbox}{title=\textbf{Advanced Note},
  colframe=GuideSlate,  colback=GuideSlate!6,  coltitle=white, colbacktitle=GuideSlate}
% In-line worked examples live in this box; each chapter's main closing worked
% example stays a numbered section so it appears in the table of contents.
\newtcolorbox{workedbox}[1][]{title=\textbf{Worked Example#1},
  colframe=GuidePlum,  colback=GuidePlum!5,   coltitle=white, colbacktitle=GuidePlum}

\newcolumntype{Y}{>{\raggedright\arraybackslash}X}
\newcolumntype{L}[1]{>{\raggedright\arraybackslash}p{#1}}

\renewcommand{\arraystretch}{1.18}
